% PhD-level technical monograph: improvements from baseline abfe578 to a newer stage but not exactly HEAD.
% Compile: pdflatex thesis_improvements.tex (twice for references) or latexmk -pdf
\documentclass[11pt,a4paper,twoside,openright]{report}

\usepackage[T1]{fontenc}
\usepackage[utf8]{inputenc}
\usepackage{lmodern}
\usepackage{microtype}
\usepackage{amsmath,amssymb,amsfonts,amsthm}
\usepackage{mathtools}
\usepackage{graphicx}
\usepackage{booktabs}
\usepackage{tabularx}
\usepackage{multirow}
\usepackage{algorithm}
\usepackage{algorithmic}
\usepackage{hyperref}
\usepackage{cleveref}
\usepackage{geometry}
\usepackage{setspace}
\usepackage{caption}
\usepackage{subcaption}
\usepackage{enumitem}
\usepackage{xcolor}
\usepackage{listings}

\geometry{margin=2.5cm}
\onehalfspacing
\hypersetup{colorlinks=true,linkcolor=blue!60!black,citecolor=green!50!black,urlcolor=blue!70!black}

\lstset{
  basicstyle=\ttfamily\small,
  breaklines=true,
  frame=single,
  columns=fullflexible,
}

\newtheorem{definition}{Definition}[chapter]
\newtheorem{proposition}{Proposition}[chapter]
\newtheorem{remark}{Remark}[chapter]

\title{%
  \textbf{Evolution of a SCION-Aware Reinforcement-Learning}\\[0.4em]
  \textbf{Path-Selection Simulator:}\\[0.2em]
  From Prototype Pipeline to a Reproducible\\[0.2em]
  Multi-Pair Evaluation Framework}
}
\author{Technical Report --- \texttt{scion-dqn-sim}\\[0.3em]
  Baseline commit: \texttt{abfe5786f094b9a11df8ba477bd35e39153ddaf2}\\
  Comparison target: \texttt{HEAD} + working tree (May 2026)}
\date{\today}

\begin{document}
\maketitle

\begin{abstract}
This document provides a doctoral-level account of every substantive improvement
developed in the \texttt{scion-dqn-sim} repository since commit
\texttt{abfe5786f094b9a11df8ba477bd35e39153ddaf2} (April 2026).
That baseline already contained a six-step evaluation pipeline, a BRITE-driven
topology generator, a beacon simulator, traffic scripts, a monolithic DQN trainer
embedded in \texttt{evaluation/03\_train\_dqn.py}, and six classical baselines.
The present system replaces fragmented research code with a \emph{library-first}
architecture under \texttt{src/}, a thin orchestration layer under
\texttt{evaluation/}, and a comprehensive pytest suite.
Improvements span (i)~build reproducibility and configuration management,
(ii)~faithful SCION control-plane and path-composition semantics,
(iii)~link-level traffic dynamics with multi-autonomous-system (AS) pair pools,
(iv)~a unified \texttt{EvaluationPathSelectionEnv} aligning training and evaluation,
(v)~a family of deep Q-network (DQN) agents including path-scoring architectures,
(vi)~corrected reward normalization and probe-cost accounting, and
(vii)~publication-grade visualization and regression testing.
The final chapter catalogs figures required to \emph{showcase} each improvement
in a dissertation or journal article.
\end{abstract}

\tableofcontents
\listoffigures
\listoftables

%==============================================================================
\chapter{Introduction}
%==============================================================================

\section{Problem statement}

Internet path selection in the SCION architecture differs fundamentally from
shortest-path routing in traditional inter-domain systems.
End hosts choose among \emph{explicit path records} assembled from up-, core-, and
down-segments discovered via beacon propagation.
Operational goals combine latency, available bandwidth, loss, hop count, and
\emph{trust} in path elements, while \emph{selective probing} imposes measurable
control-plane cost: each bandwidth or latency measurement consumes time and may
degrade user-perceived performance if overused.

Learning a probing and selection policy with reinforcement learning (RL) requires
a simulator that (a)~generates realistic AS-level topologies,
(b)~produces SCION-compliant candidate paths,
(c)~evolves link metrics under competing traffic, and
(d)~exposes a Markov decision process (MDP) whose reward matches evaluation metrics.
The baseline repository at \texttt{abfe578} demonstrated feasibility but suffered
from duplicated modules, inconsistent time indices in reward computation,
single-pair training bias, and evaluation scripts that ignored probing penalties.

\section{Scope of this document}

We compare the baseline tree (27~commits behind \texttt{HEAD} at the time of
writing, plus uncommitted architectural refactoring) along three axes:
\begin{enumerate}[label=(\roman*)]
  \item \textbf{Scientific fidelity}: SCION beaconing, path legality, traffic coupling.
  \item \textbf{Learning methodology}: state representations, algorithms, reward design.
  \item \textbf{Software engineering}: modularity, testing, reproducibility, documentation.
\end{enumerate}

Intermediate commits (chronological highlights):
\begin{itemize}
  \item \texttt{a60a713}: \texttt{uv}/\texttt{pyproject.toml} packaging;
  \item \texttt{fdea966}--\texttt{f3c9668}: topology visualization overhaul;
  \item \texttt{a3698f3}--\texttt{a7365e9}: unified YAML topology configuration;
  \item \texttt{2bc82f2}: cross-ISD non-core edge pruning in BRITE conversion;
  \item \texttt{3b20cc4}: large refactor of beaconing, traffic, evaluation env;
  \item \texttt{c814ad1}: beacon simulator correctness fixes;
  \item \texttt{77886fd}--\texttt{e708848}: simple DQN and multi-pair restoration;
  \item \texttt{6c70254}: path-scoring DQN (simple and enhanced);
  \item \texttt{de50858}: reward synchronization and training library extraction;
  \item Working tree: \texttt{src/pipeline/}, \texttt{link\_traffic\_sim.py}, removal of evaluation shims.
\end{itemize}

%==============================================================================
\chapter{Baseline System (\texttt{abfe578})}
%==============================================================================

\section{Pipeline topology}

The baseline evaluation flow was numbered but \emph{inconsistent}:
\begin{center}
\begin{tabular}{cl}
\toprule
Step & Script \\
\midrule
01 & \texttt{01\_generate\_topology.py}, \texttt{01\_setup\_dense\_topology.py} \\
02 & \texttt{02\_run\_beaconing.py}, \texttt{02\_simulate\_traffic.py} \\
03 & \texttt{03\_simulate\_traffic.py}, \texttt{03\_train\_dqn.py} \\
04 & \texttt{04\_train\_dqn.py} \\
05 & \texttt{05\_evaluate\_methods.py}, two figure scripts \\
06 & \texttt{06\_generate\_figures.py} \\
\bottomrule
\end{tabular}
\end{center}
Training lived in step~03 (\texttt{03\_train\_dqn.py}) with an inline
\texttt{DQN}, \texttt{ReplayBuffer}, and \texttt{SelectiveProbingEnvironment}
classes---approximately 300 lines of non-reusable logic per script.

\section{Reinforcement-learning stack}

Two parallel RL implementations coexisted:
\begin{itemize}
  \item \texttt{src/rl/evaluation/}: legacy Gym environment
    (\texttt{environment.py}, \texttt{dqn\_agent.py}) expecting NPZ flow tensors
    and high-dimensional concatenated features ($8$ path $\times$ max\_paths
    $+\,6$ network $+\,4$ temporal dimensions).
  \item \texttt{src/rl/dqn\_agent\_enhanced.py}: enhanced agent used sporadically
    by newer evaluation scripts.
\end{itemize}
Baselines wrapped selectors with \texttt{probing\_wrapper.py} and
\texttt{fixed\_source\_wrapper.py}, duplicating probe accounting outside the
environment.

\section{Topology and beaconing}

BRITE integration existed (\texttt{brite\_cfg\_gen.py},
\texttt{brite2scion\_converter.py}) but lacked:
\begin{itemize}
  \item a single declarative configuration file for experiments;
  \item geographic ISD placement helpers;
  \item cross-ISD pruning before peering augmentation;
  \item a dedicated SCION path builder enforcing up--core--down composition.
\end{itemize}
\texttt{beacon\_sim\_v2.py} operated on pandas node/edge tables without the
strict parent$\rightarrow$child intra-ISD constraint later enforced.

\section{Traffic simulation}

\texttt{03\_simulate\_traffic.py} produced hourly path metrics with limited
coupling between flows: background traffic and per-link aggregation were absent
or simplified relative to the current link-bottleneck model.

\section{Testing and packaging}

No \texttt{tests/} directory, no \texttt{pyproject.toml}, and dependency pinning
via \texttt{requirements.txt} only.
\texttt{src/cli.py} and \texttt{src/example\_usage.py} provided ad hoc entry points
that diverged from the evaluation pipeline.

%==============================================================================
\chapter{Infrastructure and Reproducibility}
%==============================================================================

\section{Modern Python packaging (\texttt{a60a713}, \texttt{f141559})}

\paragraph{Improvement.}
Introduction of \texttt{pyproject.toml} with Hatchling, \texttt{uv.lock}, and
\texttt{uv sync --extra dev} for deterministic environments; \texttt{setup\_brite.sh}
documents JAR construction (\texttt{Java/Brite.jar}) with the three-argument BRITE
invocation contract.

\paragraph{Scientific impact.}
Experiments become reproducible across machines; BRITE failures are isolated to
a single shell script rather than opaque subprocess errors inside Python.

\paragraph{Implementation artifacts.}
\texttt{pyproject.toml}, \texttt{uv.lock}, \texttt{setup\_brite.sh},
\texttt{AGENTS.md}, \texttt{GEMINI.md}.

%==============================================================================
\chapter{Topology Generation and Visualization}
%==============================================================================

\section{Unified YAML configuration (\texttt{a3698f3}, \texttt{a7365e9})}

\paragraph{Improvement.}
\texttt{evaluation/topology\_defaults.yaml} centralizes generator choice, BRITE
\texttt{java\_model} parameters, SCION converter options, and output flags.
\texttt{01\_generate\_topology.py} merges CLI overrides via
\texttt{--topology-config}; resolved config is dumped to
\texttt{topology/topology\_config\_resolved.yaml}.

\paragraph{Refinement (working tree).}
Logic migrated to \texttt{src/topology/eval\_config.py}
(\texttt{load\_unified\_topology\_config}, \texttt{run\_brite\_topology\_generation})
so the pipeline script is a thin CLI.

\section{BRITE-to-SCION converter enhancements (\texttt{2fe0c53}, \texttt{2bc82f2})}

\paragraph{Improvement.}
\texttt{BRITE2SCIONConverter} now supports:
\begin{itemize}
  \item configurable \texttt{prune\_cross\_isd\_noncore\_fraction} to remove unrealistic
    dense cross-ISD connectivity among non-core ASes before random peering injection;
  \item seeded extra-peering link injection (\texttt{extra\_peering\_max\_links},
    \texttt{extra\_peering\_seed});
  \item ISD assignment and core selection via \texttt{topology\_geo.py} (k-means placement,
    core rings, latency helpers);
  \item step-wise PNG exports (\texttt{step1\_vanilla\_brite.png}, \ldots) when enabled.
\end{itemize}

\paragraph{Theoretical note.}
Pruning reduces clique-like cross-ISD artifacts inherent in Waxman/Barabási--Albert
graphs, better matching SCION's hierarchical isolation domains.

\section{Topology visualization (\texttt{fdea966}, \texttt{f3c9668})}

\paragraph{Improvement.}
\texttt{src/visualization/topology\_visualizer.py} provides:
\begin{itemize}
  \item \textbf{full} dashboard: geographic map, degree histogram, ISD composition,
    link-type pie chart;
  \item \textbf{simple} geographic rendering with role-colored AS nodes;
  \item optional extras: ISD map, core-only subgraph, connectivity matrix;
  \item \texttt{topology\_stats.txt} report generation.
\end{itemize}
CLI moved to \texttt{src/visualization/topology\_cli.py} (working tree).

%==============================================================================
\chapter{SCION Control Plane and Path Discovery}
%==============================================================================

\section{Corrected beacon simulator (\texttt{2a78053}, \texttt{c814ad1}, \texttt{3b20cc4})}

\paragraph{Improvement.}
\texttt{CorrectedBeaconSimulator} enforces:
\begin{itemize}
  \item core beaconing only on \texttt{type=core} half-edges;
  \item intra-ISD propagation strictly parent$\rightarrow$child (reverse edges never
    carry PCBs upstream);
  \item PEER edges excluded from PCB flooding;
  \item per-origin segment fan-out cap (\texttt{max\_segments\_per\_origin}, overridable
    via \texttt{BEACON\_MAX\_SEGMENTS\_PER\_ORIGIN}).
\end{itemize}

\section{SCION path builder (\texttt{3b20cc4})}

\paragraph{Improvement.}
\texttt{src/simulation/path\_builder.py} enumerates paths obeying
Up $\rightarrow$ Core $\rightarrow$ Down segment assembly with optional peering
shortcuts, parameterized by explicit \texttt{core\_ases}.
Unit tests in \texttt{tests/test\_path\_builder\_scion.py} lock composition rules.

\section{Multi-pair path store (\texttt{e708848})}

\paragraph{Improvement.}
\texttt{02\_run\_beaconing.py} populates \texttt{InMemoryPathStore} for a
\texttt{pair\_pool} of AS pairs (not only the legacy single
\texttt{selected\_pair.json}).
Downstream training and evaluation iterate diverse $(src,dst)$ tuples, reducing
overfitting to one hot potato pair.

%==============================================================================
\chapter{Traffic and Link Dynamics}
%==============================================================================

\section{Link-level traffic simulation (\texttt{3b20cc4}, refined in \texttt{de50858})}

\paragraph{Improvement.}
The traffic model (now \texttt{src/simulation/link\_traffic\_sim.py}) implements:
\begin{align}
  \lambda_{f}(t) &= B \cdot \delta(h) \cdot \omega(d) \cdot \xi_f, \\
  U_e(t) &= \min\!\left(1.5,\; \frac{\sum_{f: e \in f} b_f(t)}{C_e}\right),
\end{align}
where $\delta$ is a diurnal factor, $\omega$ a weekend attenuation, $\xi_f$ pair-specific
jitter, and $C_e$ link capacity.

\paragraph{Foreground traffic.}
Each pair in \texttt{pair\_pool} receives a 28-day $\times$ 24-hour demand series.
Demand splits across the three lowest-latency paths with weights inversely proportional
to static latency (ECMP-like spread).

\paragraph{Background traffic.}
Up to $\min(200, 4|V|)$ random routable pairs per hour inject 10--80~Mbps on their
minimum-static-latency path, creating contention on shared edges.

\paragraph{Path metrics from links.}
For path $p$ at hour $t$, per-link queueing augments latency; path latency, available
bandwidth, utilization, and end-to-end loss aggregate via bottleneck and independence
assumptions:
\[
  \text{loss}_p = 1 - \prod_{e \in p} (1 - \text{loss}_e).
\]

\paragraph{Schema.}
\texttt{link\_states.pkl} stores \texttt{by\_pair} blocks
(\texttt{pair\_<src>\_<dst>/path\_<idx>}) plus legacy flat \texttt{path\_<idx>} keys
for backward compatibility.

%==============================================================================
\chapter{Unified Evaluation Environment}
%==============================================================================

\section{\texttt{EvaluationPathSelectionEnv} (\texttt{3b20cc4}, \texttt{de50858})}

\paragraph{Improvement.}
A single environment class (\texttt{src/simulation/evaluation\_env.py}) subsumes
observation construction, selective probing, action masking, and reward computation
for both training (step~04) and evaluation (step~05).

\subsection{Observation modes}

\begin{definition}[Flat observation]
Vector $\mathbf{s}^{\mathrm{flat}} \in \mathbb{R}^{5}$: normalized hour, mean latency,
loss, utilization, and bandwidth summaries for the probed path set.
\end{definition}

\begin{definition}[Scoring observation]
Dictionary with global $\mathbf{g} \in \mathbb{R}^{7}$ (five link aggregates plus
two-dimensional normalized pair embedding) and per-path matrix
$\mathbf{P} \in \mathbb{R}^{N \times 7}$ (latency, loss, hops, relative bandwidth,
utilization, static bandwidth, trust).
\end{definition}

\subsection{Probing model}

Probes incur costs
\[
  c_{\mathrm{probe}} = n_{\ell} c_{\ell} + n_{b} c_{b} + \sum_{h} c_{\mathrm{hop}}(h),
\]
with optional normalization per step (\texttt{normalize\_probe\_penalty}) so comparing
RL selective probing against baseline full probing is fair.

\subsection{Reward function}

Let $g$ be goodput (selected bandwidth divided by oracle max bandwidth at the
\emph{same} hour), $\tau$ trust, and $\pi$ normalized probe penalty:
\begin{equation}
  r = 2 \bigl(w_1 g + w_2 \tau\bigr) - 1 - w_{\mathrm{probe}} \cdot \pi.
  \label{eq:reward}
\end{equation}
Default weights: $w_1{=}0.7$, $w_2{=}0.3$, $w_{\mathrm{probe}}{=}0.05$.

\section{Reward synchronization fix (\texttt{de50858})}

\paragraph{Problem (baseline and early unified env).}
\texttt{step()} advanced the hour index \emph{before} reward used
\texttt{current\_link\_states}, causing goodput normalization against $t{+}1$
capacity while the action applied at $t$---documented in \texttt{REPORT.md}.

\paragraph{Solution.}
\texttt{compute\_reward} accepts explicit \texttt{max\_possible\_bw} captured at
decision time; \texttt{EvaluationPathSelectionEnv.step} caches oracle bandwidth
before advancing time. Training and evaluation call the same code path via
\texttt{path\_selection\_train.py} and \texttt{05\_evaluate\_methods.py}.

\paragraph{Empirical effect.}
On run \texttt{run\_20260528\_154104\_single\_pair}, corrected rewards cluster near
$0.85$--$0.91$ with standard deviations $0.01$--$0.07$ for strong methods, versus
bimodal $[-0.4, 0.95]$ behavior described in the diagnostic report for the buggy
normalization.

%==============================================================================
\chapter{Reinforcement-Learning Agents}
%==============================================================================

\section{Enhanced flat DQN (\texttt{dqn\_agent\_enhanced.py})}

Double DQN, dueling architecture, prioritized experience replay (PER), action masking
for invalid paths, and configurable MLP widths. Used by \texttt{04\_train\_dqn.py}.

\section{Simple flat DQN (\texttt{77886fd})}

\texttt{src/rl/dqn\_agent\_simple.py}: minimal two-hidden-layer MLP for ablation
against architectural complexity (\texttt{04\_train\_simple\_dqn.py}).

\section{Path-scoring DQNs (\texttt{6c70254})}

\paragraph{Motivation.}
Flat DQNs index actions $\{0,\ldots,K{-}1\}$; when $K$ varies across pairs, transfer
is poor. Path-scoring networks assign a scalar $Q(s, p_i)$ to each candidate path
feature row.

\subsection{Simple path-scoring DQN}
\texttt{SimplePathScoringDQNAgent}: shared MLP scoring head, variable $N$ paths,
invalid path masking with large negative logits.

\subsection{Enhanced path-scoring DQN}
\texttt{EnhancedPathScoringDQNAgent}: dueling decomposition
$Q = V(\mathbf{g}) + A(\mathbf{g}, \mathbf{p}_i) - \mathrm{mean}_j A(\mathbf{g}, \mathbf{p}_j)$,
Double DQN targets, PER, padding batches for variable $N$.

\section{Training library (\texttt{de50858})}

\texttt{src/rl/path\_selection\_train.py} centralizes:
\begin{itemize}
  \item \texttt{train\_flat\_dqn(run\_path, variant)} with adaptive episode counts;
  \item \texttt{train\_scoring\_dqn(run\_path, variant, ScoringHyperparams)};
  \item shared \texttt{TRAINING\_HOURS} $=$ first 14 days, \texttt{EPISODE\_LENGTH}${=}24$.
\end{itemize}
Evaluation scripts reduced from $\sim$300 lines to $\sim$25 lines each.

%==============================================================================
\chapter{Baselines and Evaluation Methodology}
%==============================================================================

\section{Classical selectors}

Six baselines in \texttt{src/baselines/}: shortest path, widest path, lowest latency,
ECMP, random, SCION default. Step~05 evaluates all methods on the same
\texttt{EvaluationPathSelectionEnv} instance.

\section{ECMP determinism (\texttt{be24c64}, working tree)}

ECMP now hashes the flow 5-tuple to a stable index among equal-hop paths, replacing
\texttt{np.random.choice} which broke reproducibility and failed unit tests.

\section{Evaluation harness (\texttt{05\_evaluate\_methods.py})}

\begin{itemize}
  \item Multi-pair evaluation: up to 32 pairs from \texttt{pair\_pool};
  \item Evaluation window: hours $14 \times 24$ through $28 \times 24$;
  \item Metrics: reward mean/std/percentiles, latency, bandwidth, probe counts and times,
    wall-clock selection time;
  \item Output: \texttt{evaluation\_results.json} consumed by step~06.
\end{itemize}

\section{Figure generation (\texttt{06\_generate\_figures.py})}

LNCS-style matplotlib defaults via \texttt{src/pipeline/figures.py}:
\begin{itemize}
  \item Figure~1: probe overhead and selection time (bar charts);
  \item Figure~2: reward distribution (box plots from mean/std summary);
  \item Figure~3: latency vs.\ bandwidth probe breakdown per selection.
\end{itemize}

%==============================================================================
\chapter{Software Architecture Refactoring}
%==============================================================================

\section{Separation of concerns (working tree, May 2026)}

\paragraph{Motivation.}
Evaluation directory contained library shims
(\texttt{path\_selection.py}, \texttt{train\_lib.py}, \texttt{\_common.py}) and
468-line traffic logic inside a script.

\paragraph{Result.}
\begin{center}
\begin{tabular}{ll}
\toprule
Layer & Responsibility \\
\midrule
\texttt{evaluation/} & Numbered CLIs, \texttt{topology\_defaults.yaml}, \texttt{\_bootstrap.py} \\
\texttt{src/pipeline/} & Run-dir resolution, subprocess runner, figure styling \\
\texttt{src/simulation/} & Path store, env, traffic sim, run context \\
\texttt{src/topology/} & BRITE, conversion, eval config \\
\texttt{src/rl/} & Agents and training loops \\
\texttt{tests/} & 58 pytest cases, no evaluation imports \\
\bottomrule
\end{tabular}
\end{center}

\paragraph{Orchestrator.}
\texttt{evaluation/run\_full\_evaluation.py} supports \texttt{--from-step} and
\texttt{--run-dir}, restoring a documented entry point missing in intermediate commits.

%==============================================================================
\chapter{Verification and Quality Assurance}
%==============================================================================

\section{Test suite}

Fifty-eight tests cover:
\begin{itemize}
  \item BRITE config generation and cross-ISD pruning;
  \item beacon simulator smoke;
  \item SCION path builder legality;
  \item \texttt{InMemoryPathStore} serialization;
  \item \texttt{EvaluationPathSelectionEnv} reward timing and observations;
  \item all six baselines;
  \item DQN and path-scoring DQN smoke;
  \item pipeline helpers (\texttt{src/pipeline/});
  \item topology unified config file presence.
\end{itemize}

\section{Diagnostic reporting}

\texttt{REPORT.md} records root-cause analysis of poor rewards (bandwidth exhaustion,
temporal desynchronization, missing probe penalties) and guided the \texttt{de50858}
fixes---included here as scientific audit trail.

%==============================================================================
\chapter{Empirical Summary}
%==============================================================================

Table~\ref{tab:sample-run} excerpts run \texttt{run\_20260528\_154104\_single\_pair}
after reward fixes (single-pair configuration; multi-pair runs may differ).

\begin{table}[htbp]
\centering
\caption{Sample evaluation summary (336 selections per method).}
\label{tab:sample-run}
\begin{tabular}{lcccc}
\toprule
Method & Reward mean & Reward std & Avg probes/sel. & Latency (ms) \\
\midrule
Scoring simple DQN & 0.913 & 0.009 & 140.1 & 22.8 \\
Scoring enhanced DQN & 0.903 & 0.016 & 142.0 & 24.8 \\
Widest path & 0.900 & 0.009 & 4640.0 & 28.2 \\
Enhanced DQN & 0.887 & 0.038 & 145.2 & 24.3 \\
Shortest path & 0.885 & 0.047 & 341.0 & 29.0 \\
Simple DQN & 0.855 & 0.075 & 155.6 & 29.5 \\
Lowest latency & 0.743 & 0.130 & 341.0 & 17.8 \\
\bottomrule
\end{tabular}
\end{table}

\begin{remark}
Widest-path achieves high reward only when probe penalties are dominated by
goodput; with full $w_{\mathrm{probe}}$ enforcement, RL methods dominate
\emph{probe--reward Pareto} fronts. Always report probe overhead alongside raw reward.
\end{remark}

%==============================================================================
\chapter{Required Figures to Showcase Improvements}
%==============================================================================

This chapter lists \textbf{every figure} recommended to substantiate the improvements
above. Existing pipeline outputs are marked \textbf{[auto]}; others require
\textbf{[new]} scripts or post-processing.

\section{Infrastructure and reproducibility}

\begin{enumerate}
  \item \textbf{[new] Dependency closure diagram}: \texttt{uv.lock} package graph
    or table of pinned versions vs.\ baseline \texttt{requirements.txt}.
  \item \textbf{[new] BRITE build verification screenshot}: successful
    \texttt{setup\_brite.sh} output and \texttt{Brite.jar} timestamp.
\end{enumerate}

\section{Topology generation and configuration}

\begin{enumerate}[resume]
  \item \textbf{[auto]} \texttt{topology/step1\_vanilla\_brite.png}: raw BRITE graph.
  \item \textbf{[auto]} \texttt{topology/step2\_scion\_classified.png}: ISD coloring and roles.
  \item \textbf{[auto]} \texttt{topology/step3\_peering\_enhanced.png}: after peering injection.
  \item \textbf{[new] Pruning ablation}: two geographic maps side-by-side,
    fraction $\in \{0, 0.2, 0.35\}$ cross-ISD non-core edges; metrics: edge count,
    average shortest path length, cross-ISD diameter.
  \item \textbf{[auto]} \texttt{topology\_dashboard.png} (full mode): publication overview.
  \item \textbf{[auto]} \texttt{topology\_geographic.png} (simple mode).
  \item \textbf{[auto]} \texttt{isd\_map.png}, \texttt{core\_network.png},
    \texttt{connectivity\_matrix.png}.
  \item \textbf{[auto]} \texttt{topology\_stats.txt} rendered as bar/table figure:
    AS count, link type histogram, core fraction.
  \item \textbf{[new]} YAML sensitivity: BRITE $N$ nodes vs.\ mean path count per pair
    (error bars over 5 seeds).
\end{enumerate}

\section{Beaconing and path discovery}

\begin{enumerate}[resume]
  \item \textbf{[new]} PCB propagation sanity: time--snapshot diagram of segment counts
    per ISD (parent$\rightarrow$child only).
  \item \textbf{[new]} Path-count CDF: number of valid SCION paths per pair before/after
    \texttt{max\_segments\_per\_origin} cap.
  \item \textbf{[auto]} \texttt{beacon\_output/beacon\_stats\_corrected.json} visualized:
    segments per origin histogram.
  \item \textbf{[new]} Multi-pair coverage heatmap: $src \times dst$ matrix of path
    availability (binary or count).
  \item \textbf{[new]} Path composition stacked bar: fraction of paths using
    peering vs.\ core-only routes.
\end{enumerate}

\section{Traffic and link dynamics}

\begin{enumerate}[resume]
  \item \textbf{[auto]} \texttt{traffic\_flows.csv}: diurnal demand curve for selected pair
    (28-day overlay or mean $\pm$ std by hour-of-day).
  \item \textbf{[new]} Link utilization time series: top-$k$ congested edges, hours 0--672.
  \item \textbf{[new]} Background vs.\ foreground load stacked area chart (total Mbps).
  \item \textbf{[new]} Bottleneck distribution: histogram of path available bandwidth
    at evaluation hours.
  \item \textbf{[new]} Queueing latency multiplier vs.\ utilization scatter per link
    (validate $q(U_e)$ model).
  \item \textbf{[new]} Multi-pair contention: CDF of pairs with zero available bandwidth
    under multi-pair vs.\ single-pair simulation.
\end{enumerate}

\section{Environment, reward, and probing}

\begin{enumerate}[resume]
  \item \textbf{[new]} Reward ablation bars: variants of \eqref{eq:reward} with
    $w_{\mathrm{probe}}{=}0$ vs.\ $0.05$, with/without oracle \texttt{max\_possible\_bw}.
  \item \textbf{[new]} Off-by-one demonstration: reward difference histogram
    (buggy $t$ vs.\ $t{+}1$ normalization) on fixed actions.
  \item \textbf{[new]} Goodput vs.\ raw bandwidth scatter colored by hour (show saturation).
  \item \textbf{[new]} Probe-cost breakdown pie per method (latency vs.\ bandwidth probes).
  \item \textbf{[auto]} \texttt{figure1\_probe\_overhead.png}: avg probe ms and selection time.
  \item \textbf{[auto]} \texttt{figure3\_probe\_breakdown.png}: probes per selection by type.
\end{enumerate}

\section{Reinforcement-learning agents}

\begin{enumerate}[resume]
  \item \textbf{[auto]} \texttt{training\_stats.json}, \texttt{dqn\_*\_training\_stats.json}:
    learning curves (episode reward vs.\ episode index) for all four DQN variants.
  \item \textbf{[new]} Sample efficiency: steps to reach 90\% of final reward vs.\ method.
  \item \textbf{[new]} Action distribution entropy over paths during training (exploration).
  \item \textbf{[new]} Path-scoring interpretability: heatmap of $Q(s,p_i)$ vs.\ path rank
    by static latency for one representative state.
  \item \textbf{[new]} Architecture ablation table: enhanced vs.\ simple vs.\ scoring
    (parameters, inference ms, reward mean).
\end{enumerate}

\section{End-to-end evaluation}

\begin{enumerate}[resume]
  \item \textbf{[auto]} \texttt{figure2\_path\_reward.png}: reward distribution by method.
  \item \textbf{[new]} \textbf{Pareto front}: mean reward vs.\ mean probe time
    (scatter, all methods); highlights RL advantage over widest-path oracle.
  \item \textbf{[new]} Latency--bandwidth tradeoff scatter per method (evaluation hours).
  \item \textbf{[new]} Per-pair reward boxplots (multi-pair runs): variance across
    \texttt{pair\_pool}.
  \item \textbf{[new]} Statistical significance: paired Wilcoxon across hours between
    scoring-simple DQN and strongest baseline.
  \item \textbf{[new]} Generalization: train on subset of pairs, evaluate on held-out pairs.
\end{enumerate}

\section{Software engineering}

\begin{enumerate}[resume]
  \item \textbf{[new]} Module dependency graph (\texttt{src/} packages) before/after refactor.
  \item \textbf{[new]} Test coverage trend: pytest count vs.\ commit index (58 tests now).
  \item \textbf{[new]} Lines-of-code ratio: evaluation scripts vs.\ \texttt{src/} library
    (demonstrates thin pipeline).
\end{enumerate}

\section{Summary matrix}

\begin{table}[htbp]
\centering
\caption{Figure priorities for a dissertation chapter.}
\label{tab:figure-priority}
\small
\begin{tabular}{@{}p{0.42\textwidth}cc@{}}
\toprule
Figure theme & Priority & Pipeline step \\
\midrule
Topology dashboard + step PNGs & High & 01 \\
Pruning / $N$-nodes ablation & High & 01 \\
Multi-pair path heatmap & High & 02 \\
Diurnal traffic + link utilization & High & 03 \\
Reward/Pareto/probe (Figs 1--3) & High & 05--06 \\
Training curves (4 DQNs) & High & 04 \\
Reward off-by-one ablation & Medium & 05 (post) \\
Beacon PCB sanity & Medium & 02 \\
Path-scoring $Q$ heatmap & Medium & 04--05 \\
Held-out pair generalization & Medium & 04--05 \\
LOC / test coverage & Low & N/A \\
\bottomrule
\end{tabular}
\end{table}

%==============================================================================
\chapter{Conclusion}
%==============================================================================

Starting from commit \texttt{abfe578}, the \texttt{scion-dqn-sim} project evolved from
a demonstrator with duplicated scripts and inconsistent RL interfaces into a
reproducible research platform.
The critical scientific advances are:
(i)~SCION-faithful beaconing and path composition,
(ii)~link-coupled traffic with multi-pair pools,
(iii)~a unified MDP with synchronized rewards and probe accounting, and
(iv)~path-scoring DQNs that generalize across variable candidate sets.
Software refactoring into \texttt{src/} ensures these advances remain maintainable
as topology sizes, probe models, and algorithms grow.

Future work should replace synthetic reward box plots in step~06 with empirical
per-selection samples, add held-out pair generalization experiments, and publish
the pruning and background-traffic parameters as formal sensitivity analyses.

%==============================================================================
\begin{thebibliography}{99}
\bibitem{scion}
P.~Szalachowski, A.~Perrig, S.~Wandert, M.~Khatsenko, and A.~Perrig,
``SCION: A Secure Internet Architecture,'' Springer, 2017.

\bibitem{mnih2015}
V.~Mnih et al.,
``Human-level control through deep reinforcement learning,''
\textit{Nature}, vol.~518, pp.~529--533, 2015.

\bibitem{wang2016}
H.~van Hasselt, A.~Guez, and D.~Silver,
``Deep Reinforcement Learning with Double Q-learning,''
\textit{AAAI}, 2016.

\bibitem{schaul2016}
T.~Schaul, J.~Quan, I.~Antonoglou, and D.~Silver,
``Prioritized Experience Replay,''
\textit{ICLR}, 2016.

\bibitem{brite}
A.~Medina, I.~Matt, and J.~Byers,
``On the Origin of Power Laws in Internet Topologies,''
\textit{ACM SIGCOMM CCR}, vol.~30, no.~4, 2000.
\end{thebibliography}

\end{document}
